\documentclass[11pt,a4paper]{article}
\usepackage[T1]{fontenc}
\usepackage[utf8]{inputenc}
\usepackage[main=italian,provide=*]{babel}
\usepackage{amsmath,amssymb}
\usepackage{geometry}
\geometry{margin=2.3cm,top=2.5cm,bottom=2.7cm}
\usepackage{booktabs}
\usepackage[table]{xcolor}
\usepackage{tcolorbox}
\tcbuselibrary{skins,breakable}
\usepackage{titlesec}
\usepackage{enumitem}
\usepackage{needspace}
\usepackage{fancyhdr}
\usepackage{hyperref}
\hypersetup{colorlinks=true,linkcolor=Sepia,citecolor=Sepia,urlcolor=Sepia}
\newcommand{\Tr}{\operatorname{Tr}}
\newcommand{\ket}[1]{\lvert #1 \rangle}
\newcommand{\bra}[1]{\langle #1 \rvert}

% --- palette ---
\definecolor{inkblue}{HTML}{1B3A5C}
\definecolor{lightblue}{HTML}{EAF1F8}
\definecolor{accent}{HTML}{B0562A}
\definecolor{softgray}{HTML}{6B6B6B}
\definecolor{okgreen}{HTML}{2E6B3E}
\definecolor{okgreenbg}{HTML}{EAF5EC}

% --- titoli di sezione ---
\titleformat{\section}
  {\normalfont\Large\bfseries\color{inkblue}}
  {\thesection}{0.6em}{}[{\color{inkblue!40}\titlerule}]
\titlespacing{\section}{0pt}{0.7em}{0.4em}
\titleformat{\subsection}
  {\normalfont\large\bfseries\color{accent}}
  {}{0em}{}
\titlespacing{\subsection}{0pt}{0.5em}{0.2em}

\pagestyle{fancy}
\fancyhf{}
\renewcommand{\headrulewidth}{0.4pt}
\fancyhead[L]{\small\color{softgray}VQE con termine DM sotto rumore --- Passo 2, Parte 2}
\fancyhead[R]{\small\color{softgray}\thepage}
\fancyfoot[C]{\small\color{softgray}Tesi triennale, Samuele --- Relatore: Prof.\ Alessandro Chiesa}

% --- box riutilizzabili ---
\newtcolorbox{keyeq}[1][]{
  colback=lightblue, colframe=inkblue!70, boxrule=0.6pt,
  arc=2pt, left=8pt, right=8pt, top=3pt, bottom=3pt, #1
}
\newtcolorbox{okbox}[1]{
  colback=okgreenbg, colframe=okgreen!60, boxrule=0.6pt,
  arc=2pt, left=7pt, right=7pt, top=2pt, bottom=2pt,
  fonttitle=\bfseries\color{okgreen}, title={#1}
}

\title{\vspace{-1.5em}\color{inkblue}VQE con termine DM sotto rumore --- Passo 2\\[0.2em]
\Large\normalfont\color{softgray}Energia e fedeltà rumorose, punto di lavoro test 2}
\author{Note di lavoro --- Tesi triennale, Samuele\\ Relatore: Prof.\ Alessandro Chiesa}
\date{}

\begin{document}
\sloppy
\maketitle
\thispagestyle{fancy}
\vspace{-2.2em}

File: \texttt{vqe\_dm\_rumoroso\_dimero.py}, \texttt{validate\_vqe\_dm\_rumoroso\_dimero.py}.
Punto di lavoro: $b/J=0.35$, $D/J=0.80$, $J=1$ (``test 2''). Modello di
rumore: \texttt{noise\_model\_dimero.py} (Passo 1).

\section*{Metodologia}
\addcontentsline{toc}{section}{Metodologia}

\begin{itemize}[leftmargin=1.3em,itemsep=0.3em]
\item \textbf{Scelta di modellazione dichiarata}: si riusano i parametri
dell'ansatz PMA-2q$\cdot$3 già ottimizzati nel caso ideale
(\texttt{ground\_state\_test2.npz}), \emph{senza} rioptimizzare sotto
rumore. Due domande diverse, da non confondere:
\begin{itemize}[leftmargin=1.3em,itemsep=0.15em]
\item \emph{degradazione di un risultato ideale} (quello che si fa qui):
i parametri $\theta^*$ sono già stati trovati su un simulatore
\textbf{senza} rumore (Parte 1); si costruisce il circuito con quella
ricetta e lo si esegue su un simulatore \textbf{con} rumore. Domanda:
quanto peggiora un risultato già buono se eseguito su hardware
imperfetto?
\item \emph{VQE noise-aware} (non fatto qui): rifare l'intero ciclo di
ottimizzazione --- iterativo, decine o centinaia di valutazioni del
circuito --- con il rumore acceso fin dall'inizio, così che
l'ottimizzatore non veda mai un'energia pulita e converga (magari) a
parametri diversi da $\theta^*$, potenzialmente più robusti al rumore
specifico presente durante la ricerca. È lo scenario di un'esecuzione
reale su hardware, dove non esiste una versione ideale a cui accedere
per calibrare l'ottimizzazione.
\end{itemize}
La prima strada è più economica (nessuna rioptimizzazione) e isola una
domanda pulita, senza mischiarla con ``l'ottimizzatore trova comunque un
buon punto anche partendo da zero col rumore acceso''. Resta
un'estensione naturale se resterà tempo. Nota a margine, non ancora
verificata: il canale depolarizzante usato qui è isotropo (restringe la
sfera di Bloch uniformemente in ogni direzione, Fig.~2 di
\texttt{teoria\_modello\_rumore.tex}) --- un rumore di questo tipo è
tipicamente poco compensabile riaggiustando i parametri, a differenza di
un errore sistematico/coerente; ci si aspetterebbe quindi che un VQE
noise-aware non cambi granché rispetto ai risultati di questo Passo, ma è
una previsione da verificare, non un fatto acquisito.
\item \textbf{Fedeltà per uno stato misto.} La fedeltà di Uhlmann fra due
operatori densità è $F(\rho,\sigma)=\big(\Tr\sqrt{\sqrt\rho\,\sigma\,\sqrt\rho}\,\big)^2$.
Quando uno dei due è puro, $\sigma=\ket\psi\bra\psi$, si semplifica: usando
$\sqrt\sigma=\sigma$ (proiettore) e $\sigma\rho\,\sigma=(\bra\psi\rho\ket\psi)\,\sigma$,
\begin{keyeq}
\centering
$F(\rho,\ket\psi\bra\psi) = \bra\psi\rho\ket\psi$
\end{keyeq}
--- niente radice quadrata di matrice da calcolare, solo un valore di
aspettazione. Verificata contro \texttt{qiskit.quantum\_info.state\_fidelity}
su uno stato misto casuale (non il caso VQE, dove l'infedeltà è troppo
piccola per essere un test stringente).
\item \textbf{Energia rumorosa}: $E=\Tr[\rho H]$, stesso calcolo già usato
nella verifica del Passo 1.
\end{itemize}

\section*{Verifica: limite di rumore nullo}
\addcontentsline{toc}{section}{Verifica: limite di rumore nullo}

\begin{okbox}{Limite di rumore nullo verificato}
\begin{center}
\renewcommand{\arraystretch}{1.25}
\begin{tabular}{@{}lc@{}}
\toprule
\textbf{Quantità} & \textbf{Valore}\\
\midrule
$F$ (registrata, Parte 1) & $0.999999999984$\\
$F$ (ricalcolata, rumore nullo) & $0.999999999984$\\
$|\Delta F|$ & $6.66\times10^{-16}$\\
$E$ (rumore nullo) & $-3.57321145$ (identico al registrato)\\
\bottomrule
\end{tabular}
\end{center}
\end{okbox}

\section*{Risultato: rumore di riferimento (\texttt{ibm\_torino})}
\addcontentsline{toc}{section}{Risultato}

\begin{center}
\renewcommand{\arraystretch}{1.25}
\begin{tabular}{@{}lc@{}}
\toprule
\textbf{Quantità} & \textbf{Valore}\\
\midrule
$E$ (rumoroso) & $-3.50164605$\\
scostamento da $E_0$ esatto & $+0.071565$\\
$F$ (rumorosa) & $0.98500601$\\
perdita di fedeltà $1-F$ & $0.014994$\\
CNOT nel circuito & $2$ (invariato rispetto al caso ideale)\\
\bottomrule
\end{tabular}
\end{center}

\noindent Circa l'1.5\% di perdita di fedeltà per un ansatz a soli 2 CNOT
--- coerente con l'ordine di grandezza di $\lambda_{2q}\approx5\times10^{-3}$
per gate coinvolto.

\section*{Scan su $\varepsilon_{2q}$}
\addcontentsline{toc}{section}{Scan}

A parità di $\varepsilon_{1q}=2.9\times10^{-4}$ e
$p_\text{readout}=2.3\times10^{-2}$: degradazione monotona sia in energia
sia in fedeltà, come atteso.

\begin{center}
\renewcommand{\arraystretch}{1.2}
\begin{tabular}{@{}ccc@{}}
\toprule
$\varepsilon_{2q}$ & $E$ & $F$ \\
\midrule
$0$ & $-3.53740090$ & $0.99251106$ \\
$1\times10^{-3}$ & $-3.52797412$ & $0.99053235$ \\
$3.8\times10^{-3}$ (riferimento) & $-3.50164605$ & $0.98500601$ \\
$8\times10^{-3}$ & $-3.46233882$ & $0.97675530$ \\
$1.5\times10^{-2}$ & $-3.39731982$ & $0.96310762$ \\
\bottomrule
\end{tabular}
\end{center}

\needspace{6\baselineskip}
\section*{Prossimo passo}
\addcontentsline{toc}{section}{Prossimo passo}

Passo 3: quantum simulation (Trotter) sotto rumore. Qui entra per la prima
volta la strategia di transpilazione per blocco decisa nella sessione
precedente (compilare una volta il singolo passo, ripeterlo $N$ volte senza
ritranspilare) --- finora non necessaria (il VQE non ha passi ripetuti),
diventa cruciale da questo punto in poi.

\end{document}
